\chapter{总结与展望}

\section{本文工作总结}

本文针对电机故障诊断中特征提取困难、小样本诊断精度低等问题，开展了基于深度学习的电机故障诊断方法研究，主要工作和创新点总结如下：

（1）提出了一种基于CNN的端到端电机故障诊断方法。设计了一维CNN网络结构，直接从原始振动信号中自动学习故障特征，避免了传统方法中复杂的人工特征提取过程，在CWRU数据集上达到了97.3\%的诊断准确率。

（2）提出了基于迁移学习的小样本电机故障诊断方法。采用预训练-微调策略，利用源域数据预训练模型，在目标域小样本条件下进行微调，在每类仅10个训练样本的条件下仍能达到85.6\%的诊断准确率，显著优于直接训练的方法。

（3）提出了基于多尺度特征融合的电机故障诊断方法。设计MSFF-Net网络结构，利用不同大小的卷积核提取多尺度特征，并通过注意力机制进行特征融合，在CWRU数据集上达到99.2\%的诊断准确率，在江南大学数据集上达到95.6\%。

\section{存在的不足}

尽管本文在基于深度学习的电机故障诊断方面取得了一定的成果，但仍存在以下不足：

（1）实验验证主要基于公开数据集，缺乏在实际工业场景中的应用验证。实际工业环境中的噪声干扰、工况变化等因素对诊断性能的影响有待进一步研究。

（2）本文方法主要针对单一故障的诊断，对于多故障耦合的情况尚未进行深入研究。

（3）深度学习模型的可解释性较差，难以满足工业应用中对诊断结果可信度的要求。

（4）模型的实时性有待优化，当前方法在嵌入式设备上的部署和推理速度需要进一步改进。

\section{未来工作展望}

针对上述不足，未来的研究工作可以从以下几个方面展开：

（1）在实际工业场景中采集数据，验证和改进所提方法的工程适用性，研究噪声鲁棒性和工况适应性。

（2）研究多故障耦合条件下的诊断方法，探索多标签分类和故障严重程度评估技术。

（3）研究可解释性深度学习方法，如注意力可视化、类激活映射等，提高诊断结果的可信度。

（4）研究模型轻量化和压缩技术，将深度学习模型部署到嵌入式设备上，实现实时在线故障诊断。

（5）结合数字孪生技术，构建电机故障诊断的虚实交互系统，实现故障预测与健康管理。